\honest{Entangled Truths, Contagious Lies}{Eliezer Yudkowsky}{2008}

I open with two passages from E. E. Smith's Lensman novels: a fully competent mind could reconstruct a universe from one fact, and anyone who studies a pebble for a hundred years will begin to perceive its truth. \nb{In \textsc{First Lensman} the first idea is at once declared impossible, since no fully competent mind ever exists.}

I am ``reasonably sure'' that a pebble does not fix the continents and people of the Earth, since other planets and other Everett branches would produce the same pebble. But it would seem to include our laws of physics, and so, if there are no truly free variables, the whole universe. \nb{The Everett branches assume the many-worlds interpretation, which the post takes as given.} Its crystals formed under gravity and tell you something about the planet's mass; its elements tell you something about how the planet formed. \nb{Both are true in outline: minerals record the pressure they formed under, and composition is how some meteorites were traced to Mars.} I imagine telling a geologist that a pebble came from a beach at Half Moon Bay and hearing ``You liar.'' I do not know how I would be caught, ``which is the point.''

A penny pulls on the Moon with an acceleration of about $4.5 \times 10^{-31}$ m/s$^2$, so every event is connected to its past, but no astronomer could read the penny's fall from the Moon. \nb{That figure fits a coin of one gram. A United States cent weighs 2.5 grams, which gives about $1.1 \times 10^{-30}$.} The connections we can notice are fewer, and form a network I call the Great Web of Causality, capital letters half in jest.

``Oh what a tangled web we weave, when first we practise to deceive,'' said Walter Scott. Not all lies spin out of control, I say, but some do, and I give a made-up dialogue in which a lie about a business trip leads to a lie about the boss. People often fail to imagine the facts that could expose them: genetics, evidence found years later, the marks of natural selection.

Not all lies are uncovered, I repeat. But a superintelligence scanning the Earth would find many, and Paul Ekman ``could probably read off a sizeable fraction of the world's lies right now, given a chance.'' \nb{The research up to 2008 found that people detect lies about 54\% of the time, and was divided on whether some experts do much better. The post gives no evidence for its claim.} The Great Web, I say, is ``very commonly underestimated.'' I give no evidence about what liars expect.

Is honesty the best policy? ``I don't know if I'd go that far.'' But honesty or silence exposes you to fewer risks you do not know you are taking.
